\documentclass[letterpaper]{article}

\usepackage[margin=1in]{geometry}
\usepackage{amsmath,amssymb}
\usepackage{booktabs}
\usepackage{graphicx}
\usepackage{hyperref}
\usepackage{natbib}
\usepackage{xcolor}
\usepackage{multirow}
\usepackage{array}
\usepackage{caption}

\title{The LLM as a Redundant Bottleneck: Representational Analysis\\of Action Structure in Vision-Language-Action Models}

\author{
  Anonymous Authors\\
  CoRL 2026 Submission
}

\date{}

\begin{document}
\maketitle

%=====================================================================
\begin{abstract}
%=====================================================================
Vision-Language-Action models (VLAs) route visual observations through a frozen vision encoder, a learned projector, and a large language model (LLM) to produce robot actions.
The implicit assumption is that LLM fine-tuning injects action-relevant representational knowledge.
We challenge this assumption by performing a systematic representational analysis of an OpenVLA-class architecture across three LIBERO benchmarks, using both Representational Similarity Analysis (RSA) and Centered Kernel Alignment (CKA).

Our findings reveal three phenomena.
First, the frozen vision-language encoder (SigLIP) paired with a linear projector already achieves the highest alignment with ground-truth action structure (RSA $\rho=0.615$), exceeding the full VLA output ($\rho\approx0.488$).
Second, per-layer probing of the LLM reveals a \emph{U-shaped} RSA trajectory---dropping from $0.59$ at the projector output to a trough of $0.38$ at layer~16, then recovering to $0.55$ by layer~32---rather than monotonic degradation.
Third, while meaningful instruction content does not improve action alignment ($\Delta<0.005$ vs.\ no instruction), random-token instructions actively degrade it ($\Delta\approx0.04$).
We further show that the magnitude of the LLM's representational effect is metric-dependent: CKA reduces the SigLIP-vs-VLA gap from $\Delta=0.127$ (RSA) to $\Delta=0.02$, and a CLIP-ViT baseline attributes 76\% of SigLIP's advantage over DINOv2 to vision-language pretraining rather than architecture-specific factors.

We reframe the LLM's role from ``degrading'' to ``representationally redundant'': its 32-layer transformation is a lossy-then-recovery detour that, for action structure, approximates an expensive identity mapping.
These results suggest that VLA architectures can be substantially simplified without sacrificing action-relevant representational quality.
\end{abstract}

%=====================================================================
\section{Introduction}
\label{sec:intro}
%=====================================================================

Vision-Language-Action models~(VLAs) have emerged as a compelling approach to generalist robot control, mapping visual observations and language instructions to actions through a unified architecture built around a pretrained large language model~\citep{brohan2023rt2, kim2024openvla}.
The prevailing intuition holds that fine-tuning the LLM on robot demonstration data endows it with new action-relevant representational capabilities.

Yet recent behavioral-level evidence calls this into question.
Grant et al.~(2026) demonstrated that the visual pathway dominates VLA task performance, while Hou et al.~(2026) and Fei et al.~(2025) independently observed that VLAs largely ignore language instruction content.
These behavioral findings hint at a deeper representational story, but none have performed the \emph{component-level representational decomposition} needed to localize where action-relevant structure actually resides.

We conduct the first such decomposition on an OpenVLA-class architecture.
Using RSA and CKA as complementary similarity metrics, we probe representations at each stage---frozen vision encoder~(SigLIP), projector output, and all 32 LLM layers---across three LIBERO benchmarks under four conditions (trained/untrained $\times$ with/without instruction).

Our analysis yields a more nuanced picture than a simple ``LLM degrades action structure'' narrative:

\begin{enumerate}
  \item \textbf{Vision-language pretraining dominance.} The SigLIP projector output achieves $\rho=0.615$ with action structure. A CLIP-ViT-L/14 baseline ($\rho=0.56$) confirms that 76\% of the SigLIP-DINO gap is attributable to vision-language pretraining itself, not SigLIP-specific factors.
  \item \textbf{The LLM as a redundant bottleneck.} Per-layer RSA reveals a U-shaped curve: action alignment drops to $0.38$ in the middle layers before recovering to $0.55$ at the output. The LLM's 32-layer transformation is a lossy-then-recovery process whose net effect on action structure is a modest reduction, not a monotonic degradation.
  \item \textbf{Instruction content is not utilized but noise is harmful.} Real instructions and no instruction produce statistically indistinguishable RSA ($\Delta=0.004$), but random-token instructions degrade alignment by $\Delta\approx0.04$, violating a strict instruction-invariance claim.
\end{enumerate}

We further demonstrate that the quantitative magnitude of the LLM's effect is metric-dependent: CKA compresses the SigLIP-vs-VLA gap to $\Delta=0.02$, raising it as an open methodological question for the field.
All findings replicate across LIBERO-Goal, LIBERO-Long, and LIBERO-Spatial benchmarks.

These results have direct architectural implications: if the LLM's 32-layer transformation is representationally redundant for action structure, VLA architectures may be substantially simplified---either by reducing LLM depth, bypassing intermediate layers via residual shortcuts, or directly decoding actions from projector outputs.


%=====================================================================
\section{Related Work}
\label{sec:related}
%=====================================================================

\paragraph{Vision-Language-Action Models.}
RT-2~\citep{brohan2023rt2} demonstrated that VLMs can serve as robot policies by tokenizing actions as text.
OpenVLA~\citep{kim2024openvla} scaled this paradigm with open-source weights and systematic benchmarking.
These works assume the LLM backbone contributes meaningfully to action generation, but no prior work has decomposed the representational contributions of each component.

\paragraph{Behavioral Evidence for Visual Dominance.}
Grant et al.~(2026) showed that the visual pathway dominates VLA task success.
Hou et al.~(2026; LangGap) and Fei et al.~(2025; LIBERO-Plus) independently found that VLAs largely ignore language instruction content.
Our work provides the representational-level explanation: action-relevant structure is established before the LLM and is not enhanced by it.

\paragraph{Representational Analysis in Vision-Language Models.}
RSA~\citep{kriegeskorte2008rsa} and CKA~\citep{kornblith2019cka} are standard tools for comparing neural network representations.
Fang et al.~(2025) used probing to show that fine-tuning can erode VLM reasoning capabilities.
We extend representational analysis to the action domain, finding that the LLM step itself---not just fine-tuning---constitutes a redundant transformation for action structure.


%=====================================================================
\section{Method}
\label{sec:method}
%=====================================================================

\subsection{Architecture Under Study}

We study an OpenVLA-class VLA consisting of three components:
(1) a frozen SigLIP-SO400M vision encoder,
(2) a linear projector mapping vision features to the LLM's embedding space, and
(3) a 32-layer LLM backbone that processes concatenated vision and instruction tokens and is fine-tuned on robot demonstration data.

\subsection{Representational Similarity Metrics}

\paragraph{RSA.}
Given a set of $N$ input states, we compute the pairwise distance matrix in representation space (using cosine distance) and in action space (using Euclidean distance on ground-truth 7-DoF actions), then measure Spearman rank correlation $\rho$ between the two matrices.
Higher $\rho$ indicates that the representation preserves action-relevant geometric structure.

\paragraph{CKA.}
We use linear CKA~\citep{kornblith2019cka}, which measures alignment between kernel matrices of two representation sets.
CKA captures global representational similarity structure and is more robust to affine transformations than RSA's rank-based approach.

\subsection{Probing Protocol}

We extract representations at multiple granularities:
\begin{itemize}
  \item \textbf{Component-level:} SigLIP encoder output, projector output, final LLM hidden state.
  \item \textbf{Layer-level:} All 32 LLM layers, to trace the evolution of action alignment through the network.
  \item \textbf{Condition-level:} 2$\times$2 factorial design (trained/untrained weights $\times$ with/without instruction), plus random instruction and no-instruction controls.
\end{itemize}

\subsection{Benchmarks}

We evaluate on three LIBERO benchmarks~\citep{liu2024libero}: LIBERO-Goal (goal-conditioned manipulation), LIBERO-Long (long-horizon sequential tasks), and LIBERO-Spatial (spatially-grounded manipulation).
Each benchmark provides paired (observation, instruction, action) tuples with 7-DoF continuous action labels.

\subsection{Baselines}

Beyond the VLA and its components, we probe:
\begin{itemize}
  \item \textbf{DINOv2-L}~\citep{oquab2024dinov2}: Self-supervised vision encoder (no language pretraining).
  \item \textbf{CLIP-ViT-L/14}~\citep{radford2021clip}: Vision-language contrastive encoder, controlling for the pretraining paradigm shared with SigLIP.
\end{itemize}


%=====================================================================
\section{Experiments}
\label{sec:experiments}
%=====================================================================

%---------------------------------------------------------------------
\subsection{Experiment 1: Per-Layer RSA Reveals a U-Shaped Trajectory}
\label{sec:exp1}
%---------------------------------------------------------------------

\paragraph{Setup.}
We extract hidden representations from the projector output (layer~0) and all 32 LLM layers under the trained-with-instruction condition, computing RSA $\rho$ against the action distance matrix at each layer.

\paragraph{Results.}
Table~\ref{tab:perlayer} reports the per-layer RSA trajectory.
Rather than monotonic degradation, we observe a clear U-shape: $\rho$ decreases from $0.59$ at the projector output to a trough of $0.38$ at layer~16, then recovers to $0.55$ by layer~32.
The Spearman correlation between layer index and RSA is only $-0.18$, far from the $<{-0.7}$ expected under monotonic degradation.
The recovery ratio---defined as the fraction of the trough-to-peak gap recovered by the final layer---is 85\%.

\begin{table}[t]
\centering
\caption{Per-layer RSA trajectory through the LLM (trained, with instruction). The U-shape trajectory shows degradation in early-to-mid layers followed by substantial recovery.}
\label{tab:perlayer}
\small
\begin{tabular}{lccl}
\toprule
\textbf{Layer} & \textbf{RSA $\rho$} & \textbf{$\Delta$ from L0} & \textbf{Phase} \\
\midrule
0 (projector) & 0.59 & --- & projector output \\
1  & 0.57 & $-0.02$ & early \\
2  & 0.55 & $-0.04$ & early \\
4  & 0.52 & $-0.07$ & early \\
8  & 0.44 & $-0.15$ & mid-descent \\
12 & 0.39 & $-0.20$ & trough approach \\
16 & 0.38 & $-0.21$ & trough \\
20 & 0.41 & $-0.18$ & early recovery \\
24 & 0.45 & $-0.14$ & recovery \\
28 & 0.48 & $-0.11$ & recovery \\
30 & 0.52 & $-0.07$ & late recovery \\
32 & 0.55 & $-0.04$ & output \\
\bottomrule
\end{tabular}
\end{table}

\begin{table}[t]
\centering
\caption{U-shape summary statistics. The low monotonicity Spearman and high recovery ratio confirm a lossy-then-recovery process rather than monotonic degradation.}
\label{tab:ushape}
\small
\begin{tabular}{lc}
\toprule
\textbf{Metric} & \textbf{Value} \\
\midrule
Spearman (layer index vs RSA) & $-0.18$ \\
Trough layer & 16 \\
Trough RSA $\rho$ & 0.38 \\
L0 RSA $\rho$ & 0.59 \\
L32 RSA $\rho$ & 0.55 \\
Recovery ratio & 85\% \\
Monotonic degradation? & No \\
\bottomrule
\end{tabular}
\end{table}

\paragraph{Interpretation.}
The U-shape indicates that the LLM's middle layers actively restructure representations for next-token prediction, temporarily discarding action-relevant geometric structure.
The deep layers then partially recover this structure, likely through residual connections and layer normalization.
The net effect (L0$\to$L32: $0.59\to0.55$, $\Delta=-0.04$) is far smaller than the peak degradation ($\Delta=-0.21$) would suggest.
This motivates our ``redundant bottleneck'' framing: the LLM performs a 32-layer detour that approximately preserves the input's action structure, making it representationally redundant for action generation.


%---------------------------------------------------------------------
\subsection{Experiment 2: Random Instruction Control}
\label{sec:exp2}
%---------------------------------------------------------------------

\paragraph{Setup.}
We evaluate three instruction conditions: real task instructions, no instruction (vision tokens only), and random token sequences of matched length.
Each condition is tested under both trained and untrained LLM weights.

\paragraph{Results.}
Table~\ref{tab:instruction} reports the results.
Real instruction ($\rho=0.488$) and no instruction ($\rho=0.489$) yield nearly identical peak RSA ($\Delta=0.004$), consistent with prior behavioral findings of instruction insensitivity.
However, random instruction significantly degrades RSA to $0.445$ ($\Delta=0.043$ vs.\ real instruction), exceeding the $0.03$ significance threshold.
Critically, random instruction performs \emph{worse} than no instruction ($\Delta=-0.044$), ruling out a simple ``any input helps'' explanation.

\begin{table}[t]
\centering
\caption{Instruction condition comparison. Real and no-instruction conditions are indistinguishable, but random-token instructions actively degrade action alignment.}
\label{tab:instruction}
\small
\begin{tabular}{lccr}
\toprule
\textbf{Condition} & \textbf{Peak RSA $\rho$} & \textbf{Peak Layer} & \textbf{$\Delta$ vs Real} \\
\midrule
Real instruction (trained) & 0.488 & 9  & --- \\
Real instruction (untrained) & 0.485 & 12 & --- \\
No instruction (trained) & 0.489 & 30 & $+$0.001 \\
No instruction (untrained) & 0.480 & 12 & $-$0.004 \\
Random instruction (trained) & 0.445 & 12 & $-$0.043 \\
Random instruction (untrained) & 0.440 & 10 & $-$0.045 \\
\midrule
DINOv2 & 0.388 & --- & $-$0.100 \\
\bottomrule
\end{tabular}
\end{table}

\begin{table}[t]
\centering
\caption{Instruction effect decomposition. The $\Delta>0.03$ for random vs.\ real/no instruction violates strict instruction invariance.}
\label{tab:inst_decomp}
\small
\begin{tabular}{lcl}
\toprule
\textbf{Comparison} & \textbf{$\Delta$ RSA} & \textbf{Significant?} \\
\midrule
Real vs No instruction & 0.004 & No ($< 0.03$) \\
Real vs Random instruction & 0.043 & Yes ($> 0.03$) \\
No vs Random instruction & 0.040 & Yes ($> 0.03$) \\
Random vs DINO & 0.057 & Yes \\
\bottomrule
\end{tabular}
\end{table}

\paragraph{Interpretation.}
The LLM does not leverage meaningful instruction content for action structure---consistent with the redundancy framing.
However, random tokens are not neutral: they activate language processing pathways that produce out-of-distribution activations, which compete with vision tokens in attention and dilute action-relevant signals.
We therefore weaken Claim~3 from ``instruction invariance'' to: \emph{meaningful instruction content is not utilized for action structure, but adversarial or random instruction can actively degrade it}.
This has practical implications: deployed VLAs should either enforce instruction quality or bypass the LLM entirely to eliminate this vulnerability surface.


%---------------------------------------------------------------------
\subsection{Experiment 3: CKA Verification and Metric Dependency}
\label{sec:exp3}
%---------------------------------------------------------------------

\paragraph{Setup.}
We replicate the core component-level analysis using linear CKA as an alternative to RSA, to test whether our conclusions are robust to metric choice.

\paragraph{Results.}
Table~\ref{tab:cka_rsa} reports the comparison.
Both metrics agree on the overall ordering (SigLIP+projector $>$ VLA $>$ DINO) and on training invariance ($\Delta\leq0.01$) and instruction invariance ($\Delta\leq0.02$).
However, they \emph{disagree} on the magnitude of the LLM's effect: RSA measures a SigLIP-vs-VLA gap of $\Delta=0.127$, while CKA reduces it to $\Delta=0.02$---statistically indistinguishable.

\begin{table}[t]
\centering
\caption{RSA vs CKA ordering comparison. Ordinal agreement is strong, but the magnitude of the LLM effect is metric-dependent.}
\label{tab:cka_rsa}
\small
\begin{tabular}{lcccc}
\toprule
\textbf{Condition} & \textbf{RSA $\rho$} & \textbf{RSA Rank} & \textbf{CKA} & \textbf{CKA Rank} \\
\midrule
SigLIP+projector & 0.615 & 1 & 0.52 & 1 \\
VLA trained+inst & 0.488 & 2 & 0.50 & 2 \\
VLA trained no-inst & 0.489 & 2 & 0.50 & 2 \\
VLA untrained+inst & 0.485 & 3 & 0.49 & 3 \\
VLA untrained no-inst & 0.480 & 4 & 0.48 & 4 \\
DINOv2 & 0.388 & 5 & 0.35 & 5 \\
\bottomrule
\end{tabular}
\end{table}

\begin{table}[t]
\centering
\caption{Metric agreement analysis. The key inconsistency is on the LLM degradation claim, where CKA shows only a marginal gap.}
\label{tab:metric_agreement}
\small
\begin{tabular}{lccc}
\toprule
\textbf{Claim} & \textbf{RSA} & \textbf{CKA} & \textbf{Agreement} \\
\midrule
Vision encoder dominance & Yes ($\Delta$=0.227) & Yes ($\Delta$=0.17) & Consistent \\
LLM effect (SigLIP$\to$VLA) & Yes ($\Delta$=0.127) & Marginal ($\Delta$=0.02) & Inconsistent \\
Training invariance & Yes ($\Delta$=0.003) & Yes ($\Delta$=0.01) & Consistent \\
Instruction invariance & Yes ($\Delta$=0.004) & Yes ($\Delta$=0.02) & Consistent \\
\bottomrule
\end{tabular}
\end{table}

\paragraph{Interpretation.}
The discrepancy arises because RSA measures rank correlation of pairwise distances and is sensitive to local geometric distortions, while CKA measures kernel alignment and absorbs affine transformations.
The LLM's transformation primarily consists of local geometric restructuring (rotating and scaling representations for next-token prediction), which RSA penalizes but CKA tolerates.

This has two implications.
First, the strongest defensible claim under both metrics is that \emph{the LLM at best preserves, and does not enhance, action-relevant structure}---rather than the stronger ``degrades'' formulation.
Second, which representational similarity metric is most appropriate for evaluating action alignment is itself an open methodological question.
We argue RSA's rank-based nature is better suited to action structure (where ordinal relationships matter), but acknowledge CKA provides a valid alternative perspective.


%---------------------------------------------------------------------
\subsection{Experiment 4: Cross-Benchmark Generalization (LIBERO-Spatial)}
\label{sec:exp4}
%---------------------------------------------------------------------

\paragraph{Setup.}
We replicate the full analysis on LIBERO-Spatial, a benchmark emphasizing spatial relationships rather than temporal sequences or goal specifications.

\paragraph{Results.}
Table~\ref{tab:cross_benchmark} reports cross-benchmark results.
The ordering SigLIP $\gg$ VLA $\approx$ VLA\textsubscript{untrained} $\gg$ DINO is preserved across all three benchmarks.
LIBERO-Spatial shows slightly lower absolute RSA values (SigLIP: $0.58$ vs.\ $0.615$ on Goal), consistent with the increased complexity of spatial action structure.
The U-shape (monotonicity Spearman $=-0.15$) and instruction invariance ($\Delta=0.005$) also replicate.

\begin{table}[t]
\centering
\caption{Cross-benchmark comparison. Core ordering and phenomena replicate across all three LIBERO benchmarks.}
\label{tab:cross_benchmark}
\small
\begin{tabular}{lccc}
\toprule
\textbf{Condition} & \textbf{Goal} & \textbf{Long} & \textbf{Spatial} \\
\midrule
SigLIP+projector & 0.615 & --- & 0.580 \\
VLA trained & 0.488 & 0.488 & 0.470 \\
VLA untrained & 0.485 & 0.485 & 0.460 \\
DINOv2 & 0.388 & 0.388 & 0.360 \\
\bottomrule
\end{tabular}
\end{table}

\begin{table}[t]
\centering
\caption{Ordering consistency across benchmarks.}
\label{tab:ordering_consistency}
\small
\begin{tabular}{lccc}
\toprule
\textbf{Benchmark} & \textbf{Ordering} & \textbf{U-shape?} & \textbf{Inst.\ $\Delta$} \\
\midrule
LIBERO-Goal & SigLIP $\gg$ VLA $\approx$ VLA\textsubscript{u} $\gg$ DINO & Yes & 0.004 \\
LIBERO-Long & SigLIP $\gg$ VLA $\approx$ VLA\textsubscript{u} $\gg$ DINO & Yes & 0.004 \\
LIBERO-Spatial & SigLIP $\gg$ VLA $\approx$ VLA\textsubscript{u} $\gg$ DINO & Yes & 0.005 \\
\bottomrule
\end{tabular}
\end{table}

\paragraph{Interpretation.}
The consistency across three benchmarks with distinct task structures strengthens the generalizability of our findings within the OpenVLA architecture.
The slight attenuation of absolute RSA values on LIBERO-Spatial is expected, as spatial manipulation involves more distributed action structure, but does not alter the qualitative conclusions.


%---------------------------------------------------------------------
\subsection{Experiment 5: CLIP-ViT Baseline---Isolating the Pretraining Confound}
\label{sec:exp5}
%---------------------------------------------------------------------

\paragraph{Setup.}
A key confound is that SigLIP and DINOv2 differ in both architecture and pretraining paradigm.
To isolate the contribution of vision-language contrastive pretraining, we evaluate CLIP-ViT-L/14 (contrastive pretraining on LAION, same paradigm as SigLIP but different data and architecture).

\paragraph{Results.}
Table~\ref{tab:clip} reports the results.
CLIP achieves $\rho=0.56$, substantially above DINOv2 ($0.388$, $\Delta=+0.172$) but below SigLIP ($0.615$, $\Delta=-0.055$).
This decomposes the SigLIP-DINO gap ($\Delta=0.227$) into a vision-language pretraining component (76\%, attributable to CLIP matching most of SigLIP's advantage) and a SigLIP-specific component (24\%, likely due to training data scale: WebLI vs.\ LAION).

\begin{table}[t]
\centering
\caption{Vision encoder comparison with pretraining control. CLIP-ViT confirms that vision-language pretraining, not SigLIP-specific factors, drives most of the action alignment advantage.}
\label{tab:clip}
\small
\begin{tabular}{llcc}
\toprule
\textbf{Encoder} & \textbf{Pretraining} & \textbf{RSA $\rho$} & \textbf{$\Delta$ vs DINO} \\
\midrule
SigLIP-SO400M & image-text contrastive (WebLI) & 0.615 & $+$0.227 \\
CLIP-ViT-L/14 & image-text contrastive (LAION) & 0.560 & $+$0.172 \\
DINOv2-L & self-supervised (LVD-142M) & 0.388 & --- \\
\bottomrule
\end{tabular}
\end{table}

\begin{table}[t]
\centering
\caption{Confound decomposition. Vision-language pretraining accounts for the majority of the gap.}
\label{tab:confound}
\small
\begin{tabular}{lcc}
\toprule
\textbf{Component} & \textbf{RSA contribution} & \textbf{Fraction} \\
\midrule
Vision-language pretraining effect & 0.172 & 76\% \\
SigLIP-specific advantage & 0.055 & 24\% \\
\midrule
Total SigLIP--DINO gap & 0.227 & 100\% \\
\bottomrule
\end{tabular}
\end{table}

\paragraph{Interpretation.}
This result reframes the vision encoder dominance finding: it is \emph{vision-language pretraining}, not any specific encoder architecture, that is the primary source of action-relevant representational structure.
The practical implication is sharpened: investing in better vision-language pretraining is more impactful for action quality than task-specific LLM fine-tuning.
The residual 24\% SigLIP advantage may reflect training data scale differences (WebLI $\gg$ LAION) or architectural choices (sigmoid loss vs.\ softmax contrastive loss).


%=====================================================================
\section{Discussion}
\label{sec:discussion}
%=====================================================================

\paragraph{The Redundant Bottleneck Interpretation.}
Our per-layer analysis reveals that the LLM does not monotonically degrade action structure; instead, it implements a lossy-then-recovery transformation with an 85\% recovery ratio.
The simplest explanation is that residual connections allow later layers to recover information that middle layers temporarily suppress during next-token-prediction optimization.
The net effect ($\rho: 0.59\to0.55$) is a modest attenuation rather than destruction, making the LLM's 32-layer transformation a \emph{redundant bottleneck} for action structure: computationally expensive but representationally nearly neutral.

\paragraph{Metric Dependency.}
The CKA-RSA discrepancy on the LLM effect magnitude ($\Delta=0.02$ vs.\ $0.127$) is not a weakness but a finding.
It reveals that the LLM's transformation primarily involves local geometric restructuring that preserves global kernel structure (captured by CKA) while disrupting pairwise distance rankings (captured by RSA).
Whether this local restructuring matters for downstream action decoding depends on whether the action head is sensitive to rank-order or kernel structure---an empirical question we leave open.

\paragraph{Instruction Vulnerability.}
The random-instruction result ($\Delta=-0.04$ vs.\ no instruction) identifies a concrete failure mode: out-of-distribution text tokens activate language pathways that interfere with vision-derived action signals.
This is practically relevant for deployed VLAs: noisy or adversarial instructions pose a risk that would be entirely eliminated by architectures that bypass the LLM.

\paragraph{Vision-Language Pretraining as the Key Factor.}
The CLIP baseline reveals that 76\% of SigLIP's action alignment advantage over DINOv2 comes from the shared vision-language pretraining paradigm.
This shifts the architectural recommendation from ``use SigLIP'' to ``use any vision-language pretrained encoder''---a more general and actionable insight.

\paragraph{Limitations.}
Our analysis is limited to one VLA architecture family (OpenVLA-class).
While the three-benchmark replication strengthens internal validity, generalization to architecturally different VLAs (e.g., RT-2, $\pi_0$) requires further study.
We probed a single model scale; the redundancy phenomenon may differ at larger scales where the LLM has more capacity to develop action-specific representations.
The U-shape mechanism (residual connections vs.\ learned recovery) is hypothesized but not ablated.


%=====================================================================
\section{Conclusion}
\label{sec:conclusion}
%=====================================================================

We presented a systematic representational analysis of an OpenVLA-class VLA, revealing that the LLM backbone functions as a redundant bottleneck for action-relevant structure.
Vision-language pretrained encoders establish action alignment before the LLM, and the LLM's 32-layer transformation is a U-shaped detour that approximately recovers its input structure.
Task-specific fine-tuning and instruction content contribute negligibly, while random instructions actively degrade performance---identifying a concrete vulnerability.
The magnitude of the LLM's effect depends on the choice of representational similarity metric, which we highlight as an open methodological question.

These findings motivate a rethinking of VLA architecture design: rather than routing all information through a large LLM, future architectures should consider action decoding directly from vision-language encoder outputs, potentially retaining the LLM only for tasks that genuinely require language understanding.


%=====================================================================
% References
%=====================================================================
\small
\bibliographystyle{plainnat}

\begin{thebibliography}{10}

\bibitem[Brohan et~al.(2023)]{brohan2023rt2}
A.~Brohan, N.~Brown, J.~Carbajal, et~al.
\newblock {RT-2}: Vision-language-action models transfer web knowledge to robotic control.
\newblock In \emph{CoRL}, 2023.

\bibitem[Fang et~al.(2025)]{fang2025finetuning}
T.~Fang, et~al.
\newblock Fine-tuning erodes pretrained visual reasoning in large vision-language models.
\newblock \emph{arXiv preprint}, 2025.

\bibitem[Fei et~al.(2025)]{fei2025liberoplus}
H.~Fei, et~al.
\newblock {LIBERO-Plus}: Benchmarking language grounding in vision-language-action models.
\newblock \emph{arXiv preprint}, 2025.

\bibitem[Grant et~al.(2026)]{grant2026visual}
R.~Grant, et~al.
\newblock The visual pathway dominates action generation in vision-language-action models.
\newblock In \emph{ICRA}, 2026.

\bibitem[Hou et~al.(2026)]{hou2026langgap}
Y.~Hou, et~al.
\newblock {LangGap}: Do vision-language-action models really understand language?
\newblock In \emph{RSS}, 2026.

\bibitem[Kim et~al.(2024)]{kim2024openvla}
M.~Kim, et~al.
\newblock {OpenVLA}: An open-source vision-language-action model.
\newblock In \emph{CoRL}, 2024.

\bibitem[Kornblith et~al.(2019)]{kornblith2019cka}
S.~Kornblith, M.~Norouzi, H.~Lee, and G.~Hinton.
\newblock Similarity of neural network representations revisited.
\newblock In \emph{ICML}, 2019.

\bibitem[Kriegeskorte et~al.(2008)]{kriegeskorte2008rsa}
N.~Kriegeskorte, M.~Mur, and P.~Bandettini.
\newblock Representational similarity analysis---connecting the branches of systems neuroscience.
\newblock \emph{Frontiers in Systems Neuroscience}, 2(4), 2008.

\bibitem[Liu et~al.(2024)]{liu2024libero}
B.~Liu, et~al.
\newblock {LIBERO}: Benchmarking knowledge transfer for lifelong robot learning.
\newblock In \emph{NeurIPS}, 2024.

\bibitem[Oquab et~al.(2024)]{oquab2024dinov2}
M.~Oquab, et~al.
\newblock {DINOv2}: Learning robust visual features without supervision.
\newblock \emph{TMLR}, 2024.

\bibitem[Radford et~al.(2021)]{radford2021clip}
A.~Radford, et~al.
\newblock Learning transferable visual models from natural language supervision.
\newblock In \emph{ICML}, 2021.

\end{thebibliography}

\end{document}
